\section{RQ2: Methods at 16-bit}
\label{sec:rq2}

\begin{table*}[t]
\centering
\small
\begin{tabular}{lcccccc}
\toprule
 & \multicolumn{2}{c}{ASCEND} & \multicolumn{2}{c}{NSC Same Room} & \multicolumn{2}{c}{NSC Different Room} \\
\cmidrule(lr){2-3}\cmidrule(lr){4-5}\cmidrule(lr){6-7}
System & MER & \Emb & MER & \Emb & MER & \Emb \\
\midrule
small, off-the-shelf & 26.3 & 56.8 & 55.1 & 76.7 & 48.8 & 72.7 \\
tiny, FT & 26.1 & 53.2 & 34.5 & 36.2 & 26.8 & 29.8 \\
base, FT & 21.3 & 55.3 & 28.4 & 31.0 & 22.4 & 24.0 \\
small, FT & 14.8 & 30.6 & 21.4 & 23.8 & 19.0 & 22.9 \\
base, FT+L & 20.3 & 39.2 & 26.1 & 27.4 & 21.1 & 22.5 \\
small, FT+L & \textbf{14.6} & \textbf{29.8} & \textbf{20.6} & \textbf{21.3} & \textbf{17.1} & \textbf{18.7} \\
\bottomrule
\end{tabular}
\caption{Test error (\%) at 16-bit, greedy decoding. \Emb: embedded-word error. FT = ordinary fine-tuning; FT+L = fine-tuning with $\mathcal{L}_{\text{role}} + \lambda\mathcal{L}_{\text{sem}}$. Speaker-cluster 95\% CIs for small, FT+L: NSC Same Room MER [18.6, 23.4], Different Room [15.7, 18.6].}
\label{tab:main}
\end{table*}

\begin{table*}[t]
\centering
\small
\resizebox{\textwidth}{!}{%
\begin{tabular}{p{2mm}lcccccc}
\toprule
 & & \multicolumn{2}{c}{ASCEND} & \multicolumn{2}{c}{NSC Same Room} & \multicolumn{2}{c}{NSC Different Room} \\
\cmidrule(lr){3-4}\cmidrule(lr){5-6}\cmidrule(lr){7-8}
\multicolumn{2}{l}{System} & MER & \Emb & MER & \Emb & MER & \Emb \\
\midrule
\multicolumn{8}{l}{\emph{(a) Our losses vs.\ ordinary fine-tuning (FT) of the same size}} \\
 & base, role-weighted only & $-$0.7 & $-$14.6 & $-$0.6 & \textbf{$-$2.7} & $-$0.6 & \textbf{$-$0.9} \\
 & base, semantic only & $-$1.2 & $-$14.6 & \textbf{$-$1.3} & $-$1.7 & \textbf{$-$1.1} & $-$0.4 \\
 & base, FT+L (both) & $-$1.0 & $-$16.1 & \textbf{$-$2.3} & \textbf{$-$3.6} & \textbf{$-$1.3} & \textbf{$-$1.4} \\
 & small, role-weighted only & \textbf{$-$0.5} & $-$2.3 & +0.1 & $-$1.1 & $-$0.9 & $-$3.3 \\
 & small, FT+L (both) & $-$0.1 & $-$0.8 & \textbf{$-$0.7} & \textbf{$-$2.4} & \textbf{$-$1.9} & \textbf{$-$4.2} \\
\midrule
\multicolumn{8}{l}{\emph{(b) On-device gated re-ranking vs.\ the same model's greedy output}} \\
 & base, FT+L & \textbf{$-$0.7} & \textbf{$-$5.4} & \textbf{$-$0.4} & $-$0.4 & \textbf{$-$0.3} & $-$0.7 \\
 & small, FT+L & \textbf{$-$0.5} & \textbf{$-$3.9} & $-$0.1 & +0.2 & +0.1 & $-$0.1 \\
\bottomrule
\end{tabular}}
\caption{Paired differences in points (negative = better). The two panels use different baselines. Bold: 95\% CI of the difference excludes 0.}
\label{tab:deltas}
\end{table*}

\textbf{Answer: yes, but the gains are small; they are reliable on NSC and within noise on ASCEND (H2 partly supported).}

\paragraph{Fine-tuning matters most.}
Fine-tuning on ASCEND and NSC brings by far the largest gain, most of all on NSC, where off-the-shelf small has 55.1\% MER and small FT 21.4\% (Table~\ref{tab:main}).

\paragraph{Our losses help on NSC.}
Both losses lower embedded error, and together they work best (Table~\ref{tab:deltas}a). FT+L is significantly better than FT on both NSC test sets, for base and small, in both MER and embedded error, and its critical-word error does not rise. These gains (0.7--2.3 MER points) are larger than the 0.2 points by which a retrain of the same system moved. On ASCEND the changes are mostly within noise: there are only two test speakers, and a few long runs of repeated text dominate the embedded error. On our phone recordings, FT+L and FT are within noise of each other (small, embedded: $-$3.6 [$-$10.0, 2.7]). A direct penalty on transliterations did not help (Appendix~\ref{app:tl}).

\paragraph{Re-ranking helps only when gated.}
Re-ranking every clip lowers embedded error on Mandarin-main ASCEND (base FT+L, on-device: $-4.3$ points) but raises MER on English-main NSC (base FT+L, Different Room: $+0.5$; small FT+L, Same Room: $+0.8$). Gating fixes this (Table~\ref{tab:deltas}b). Only the 25--28\% of clips where the language-ID head hears a switch are re-ranked. The result is never significantly worse than greedy decoding, and for base FT+L it is significantly better on all three test sets in MER. Without the repetition filter, re-ranking fails badly, because candidates that repeat themselves win the length-normalised score. The head labels 95--98\% of speech frames correctly (Appendix~\ref{app:lid}).
